\chapter{機械学習による予測}
\label{ch:ml}

\chapref{ch:model} で予測のモデルの種類を説明した。この章では、それらを資産の予測に当てはめた主な研究が、何を示したかをまとめる。

\section{比較の基準になった研究}

\begin{paperbox}[Gu・Kelly・Xiu（2020）「機械学習による実証的資産価格」\cite{gkx2020}]
米国株（1957〜2016 年、約 3 万銘柄）の月次のリターンを、94 の銘柄の特性、8 つのマクロ変数、業種、それらの相互作用（合わせて約 900 の説明変数）から予測し、線形回帰、正則化回帰（エラスティックネット）、主成分回帰・PLS、一般化線形モデル、ランダムフォレスト、勾配ブースティング、1〜5 層のニューラルネットを比べた。
\begin{itemize}
\item 標本外の月次の決定係数は、どの方法でも 1\% に満たない。最も良いもの（3 層程度のニューラルネット）で 0.4\% 程度。すべての変数をそのまま入れた最小二乗は、大きく負になる（過学習）。
\item 木のモデルとニューラルネットが最も良く、予測で並べたロング・ショートのポートフォリオは、文献の代表的な回帰の戦略の成績を、場合によっては 2 倍にする。
\item どの方法も、重要な説明変数として、ほぼ同じもの（モメンタムの変種、流動性、変動率）を選ぶ。機械学習の利点は、新しい変数の発見ではなく、既知の変数の非線形な関係と相互作用を捉えることにある。
\item 層の深いネットワークは、浅いものより良くはならない。信号が弱く、データが少ないためである。
\end{itemize}
\end{paperbox}

この論文は、資産の予測の機械学習の研究の基準になった。同時に、「予測できる部分は極めて小さい」ことを、改めて数字で示した。

\section{たくさんの特性をどうまとめるか}

\begin{itemize}
\item \textbf{Kozak・Nagel・Santosh（2020）}\cite{kns2020}：多数の特性のポートフォリオから確率的割引ファクター（すべての資産の価格を説明する一つの量）を推定するとき、少数の特性だけを選ぶ（ラッソのような）疎なモデルではうまくいかず、多くの特性に少しずつ重みをつけ、主成分の分散の小さい方向を強く縮める（リッジに近い）推定が、標本外で最も良いことを示した。ただし、特性を主成分に回転させれば、少数の主成分で十分に説明できる。
\item \textbf{Freyberger・Neuhierl・Weber（2020）}\cite{fnw2020}：特性とリターンの非線形な関係を、ノンパラメトリックな方法（適応的なグループ・ラッソ）で推定し、62 の特性のうち、独立な情報を持つのは十数個にすぎないことを示した。
\item \textbf{Kelly・Pruitt・Su（2019）}\cite{kps2019}：特性を、銘柄のファクターへの感応度（ベータ）を決める変数として使う IPCA を提案した。少数（5 つ程度）の潜在的なファクターで、多くのアノマリーのアルファが説明できる。特性がリターンを予測するのは、それがリスクへの感応度を表すからだ、という解釈である（「特性は共分散である」）。
\item \textbf{Chen・Pelger・Zhu（2024）}\cite{cpz2024}：裁定がない条件を課した深層学習（敵対的な学習）で確率的割引ファクターを推定し、標本外で高いシャープレシオを得た。
\end{itemize}

\section{機械学習の成績は取れるのか}

\begin{paperbox}[Avramov・Cheng・Metzker（2023）\cite{acm2023}]
機械学習の予測による戦略の収益が、裁定の難しい銘柄（極めて小さな株、格付けの低い会社、苦境の会社、流動性の低い株）に集中していることを示した。そうした銘柄を除いたり、経済的な制約（売買のコスト、空売りの難しさ）を課したりすると、収益は大きく下がる。また、機械学習の戦略は回転率が高く、コストの影響が大きい。
\end{paperbox}

Leippold・Wang・Zhou（2022）\cite{lwz2022} は、中国の株式市場で同じ比較をし、機械学習の予測力は米国より強いこと、それが個人投資家の多い市場の性質（小型株、短期の反転）と関係することを示した。どこで予測できるかは、市場の参加者の構成に依存する。

\section{複雑さは良いことか}

Kelly・Malamud・Zhou（2024）\cite{kmz2024} は、説明変数の数が標本の数より多い「過剰にパラメータを持つ」モデル（ランダムな特徴量を極めて多く作り、ほぼ罰則なしで当てはめる）が、標本外で、単純なモデルより良い市場のタイミングの予測をすることを、理論と実証で示し、「複雑さの美徳」と呼んだ。古典的な統計の直感（パラメータは少ないほど良い）に反する結果で、議論を呼んでいる。

\section{機械学習は金融を学べるか}

Israel・Kelly・Moskowitz（2020）\cite{ikm2020} は、運用会社と研究者の立場から、金融の予測に機械学習を使う難しさを整理した。

\begin{enumerate}
\item \textbf{データが小さい}：月次のデータは数百か月しかない。画像や言語のデータの何万分の一である。
\item \textbf{信号が弱い}：予測できる部分は雑音のごく一部。
\item \textbf{市場が変わる}：予測が使われると、関係そのものが消える。
\item \textbf{解釈が要る}：顧客とリスク管理のために、なぜその売買をするのかを説明する必要がある。
\end{enumerate}

結論は、機械学習は役に立つが、その効果は、ほかの分野で見られるような劇的なものではなく、既存のクオンツの方法の「改良」の程度だ、というものである。

\section{価格の図を読む CNN}

\begin{paperbox}[Jiang・Kelly・Xiu（2023）「価格のトレンドを（再）画像化する」\cite{jkx2023}]
過去 5〜60 日の株価と売買高を、ローソク足の図のような画像にし、画像認識に使う畳み込みニューラルネット（CNN）に、その後のリターンの上げ下げを予測させた。CNN の予測は、既知のモメンタムや反転の変数を超える予測力を持ち、学習しなかった期間・国・時間の長さでも同じように機能した。CNN が見つけたパターンは、人間のテクニカル分析の図形とは必ずしも一致しない。
\end{paperbox}

\section{確立していること}

\begin{keybox}[まとめ]
\begin{description}[leftmargin=1em,style=nextline]
\item[確立している]
  標本外の予測の決定係数は月次で 1\% に満たない。非線形のモデル（木、浅いニューラルネット）は線形のモデルより少し良い。重要な変数は、モメンタム・流動性・変動率などの既知のもの。予測の収益は、取引の難しい小型株に偏り、コストで大きく減る。
\item[議論が続いている]
  深い・大きなモデルが本当に良いか（複雑さの美徳）。機械学習の戦略の収益が、公表と普及の後も残るか。
\end{description}
\end{keybox}

\begin{projbox}
LLM-BLM で、ChronoGPT のベクトル（数千次元）をリッジで反応に結びつけたのは、Kozak ほか（2020）の「多くの説明変数に少しずつ重みをつけて縮める」考え方に近い。また、Avramov ほか（2023）の結果は、小型株の S1 で予測の力が見つかっても、それが取れるかどうかは別の問題であることを示している。
\end{projbox}
